\scanpage{72}
% archive foliation: 71
\uline{Applications} Toutes les constructions (``universelles'') dans $\unsure{\unsure{\mathcal P|P}_*|\unsure{\mathcal P|P}^*}$ \deleted{\unsure{qui}} (\unsure{$Q_*'$|$Q_*$}, $\unsure{\mathcal P|P}_{0*}$, $K_*$) qui se font \unsure{précisément} en termes de sa structure $k$-lin. et sa $\otimes$-structure, \deleted{\unsure{donnent}} \inserted{se reflètent en} des constructions analogues sur un $C^\bullet\in\operatorname{Ob}B^\bullet$ muni de structure multiplicative.

\uline{Ex. 1} \deleted{$H^*(C^\bullet)$ On en \unsure{tire}} \inserted{Comme $C_*^\bullet\otimes C_*^\bullet$ est} une résolution $\Sigma_*$-inj. de $\Psi_*^0$ dans $\unsure{\mathcal P|P}_{0*}$, et comme $C_*^\bullet$ \unsure{lui-m\^eme}, on trouve
\diagram{\begin{tikzcd}
\operatorname{Tot}^\bullet(C_*^\bullet*C_*^\bullet) \arrow[r] \arrow[d] & C_*^\bullet\quad\text{défini à homotopie}\\
\unsure{\text{complexe associé}}&&
\end{tikzcd}}
près, d'où un hom
\diagram{\begin{tikzcd}
F_{C^\bullet}(C_*^\bullet*C_*^\bullet) & F(C_*^\bullet) \arrow[l] \arrow[d,equal] & \text{défini à hom. près} \\
F_{C^\bullet}(C^\bullet)\otimes_B F_{C^\bullet}(C^\bullet) \arrow[u] \arrow[d,equal] & C^\bullet & \\
C^\bullet\otimes C^\bullet &&
\end{tikzcd}}
d'où en composant un hom défini à homotopie près
\[
C^\bullet\otimes_B C^\bullet\longrightarrow C^\bullet
\]
avec conditions d'associativité, de comm. (attention aux \uline{signes} !) et d'unité --- le tout à homotopie près. \unsure{Il} en résulte une structure d'``$\unsure{\otimes B|B}$-Alg. graduée'' sur $H^*(C^\bullet)$, avec les propriétés habituelles.

